\documentclass[conference]{IEEEtran}
\IEEEoverridecommandlockouts

% --- Packages ---
\usepackage{amsmath,amssymb,amsthm,mathtools,bm}
\usepackage{cite}
\usepackage{enumitem}
\usepackage{graphicx}
\usepackage{booktabs}
\usepackage{array}
\usepackage{multirow}
\usepackage{textcomp}
\usepackage{xcolor}
\usepackage[hidelinks]{hyperref}
\usepackage{url}

% --- Theorem environments (IEEEtran-friendly, no AMS clash) ---
\newtheorem{theorem}{Theorem}
\newtheorem{proposition}[theorem]{Proposition}
\newtheorem{lemma}[theorem]{Lemma}
\newtheorem{corollary}[theorem]{Corollary}
\theoremstyle{definition}
\newtheorem{definition}[theorem]{Definition}
\newtheorem{remark}[theorem]{Remark}

% --- Notation macros ---
\newcommand{\R}{\mathbb{R}}
\newcommand{\E}{\mathbb{E}}
\newcommand{\LCE}{\mathcal{L}_{\mathrm{CE}}}
\newcommand{\LFL}{\mathcal{L}_{\mathrm{FL}}}
\newcommand{\LGFL}{\mathcal{L}_{\mathrm{GFL}}}
\newcommand{\LGFLq}{\mathcal{L}_{\mathrm{GFL}}^{(q)}}
\newcommand{\IF}{\mathrm{IF}}
\newcommand{\BD}{D_{\Psi}}

\begin{document}

\title{The Generalized Focal Loss: A Saturating Two-Parameter Family with Bounded Influence for Network Intrusion Detection}

\author{
\IEEEauthorblockN{Shahmeer Ali, Ali Rauf, Mohammad Soban}
\IEEEauthorblockA{
Department of Artificial Intelligence\\
FAST--NUCES Islamabad\\
\{i230119, i230076, i230056\}@isb.nu.edu.pk}
}

\maketitle

% =====================================================================
\begin{abstract}
We analyze the focal loss of Lin~\textit{et al.} as an M-estimator and identify two structural defects: an unbounded influence function as the predicted confidence $p_t \to 0^+$, and a homogeneous-modulator coupling that ties easy-example suppression to hard-example amplification through a single scalar $\gamma$. We propose a two-parameter generalization, the \emph{Generalized Focal Loss} (GFL), defined by replacing the polynomial modulator $(1-p_t)^\gamma$ with the Hill-type saturating kernel $\phi_{\gamma,\tau}(p_t) = (1-p_t)^\gamma / [1+\tau(1-p_t)^\gamma]$. We prove that (i) GFL recovers focal loss exactly at $\tau \to 0$ and cross-entropy at $\gamma \to 0$; (ii) the modulator is uniformly bounded above by $1/(1+\tau)$ on $[0,1]$; (iii) the per-sample logit gradient at the noise boundary is attenuated relative to focal loss by exactly $1/(1+\tau)$; (iv) a $q$-deformed extension achieves uniformly bounded influence on the closed interval $[0,1]$. We validate the bounded-influence theorem by direct gradient-norm measurement: on synthetic logits the predicted attenuation factor matches observation to $3 \times 10^{-4}\%$ relative error across $\tau \in [0, 5]$; on real CICIDS2017 samples produced by a CE-warmup baseline the prediction matches observation to within $10\%$ relative error. Downstream F1 on CICIDS2017 with the bounded-influence loss family is dataset-limited: the moderate-rare class tier (training counts $400$--$5000$) shows the predicted noise-as-regularization signature for the $q$-deformed variant, with mean F1 increasing by $+0.106$ from clean to $20\%$ BENIGN-to-minority label noise, while standard CE, FL, and GFL show degradation slopes within $\pm 0.013$ on the same tier. We provide the construction, proofs, validation methodology, and CICIDS2017 evaluation, with honest reporting of where the F1 metric does and does not reflect the predicted gradient-level effect.
\end{abstract}

% --- IEEE keywords ---
\begin{IEEEkeywords}
focal loss, label noise, bounded influence, intrusion detection, CICIDS2017, imbalanced classification, robust learning
\end{IEEEkeywords}

% =====================================================================
\section{Introduction}
\label{sec:intro}

Deep-learning network intrusion detection systems on CICIDS2017 \cite{sharafaldin2018toward} routinely report aggregate accuracy above 99\%, yet the operationally important attack classes remain hard to detect. Heartbleed appears 11 times in 2.8M flows; Infiltration appears 36 times. A model that flags every benign flow correctly while missing the 11 Heartbleed records achieves the same headline accuracy as one that catches all 11---and the difference between those two systems is the entire reason the task exists.

The dominant technique for addressing this imbalance has been data-level rebalancing: SMOTE \cite{chawla2002smote} synthesizes minority samples before training. SMOTE is effective in moderate-imbalance regimes~\cite{he2009learning} but provides no signal on the extreme tail where there are essentially no samples to interpolate between. The loss-level alternative is focal loss \cite{lin2017focal}, originally proposed for object detection, which down-weights well-classified examples through the polynomial modulator $(1-p_t)^\gamma$. Adoption in intrusion detection has been steady, but its reliability has been mixed; recent surveys report focal loss helping on some attack types and degrading on others without a clear pattern \cite{cherukuri2025deep, zhang2025deep}.

This paper begins by treating focal loss as an M-estimator and asking what its functional form actually does. Two structural defects emerge from the asymptotic analysis (Section~\ref{sec:fl-defects}). First, the loss is \emph{unbounded in influence}: as $p_t \to 0^+$ the per-sample gradient diverges as $\Theta(1/p_t)$, so a single mislabeled rare-class flow can dominate a minibatch's gradient. Second, the modulator is \emph{globally homogeneous}: the single scalar $\gamma$ controls both how aggressively easy examples are suppressed and how strongly hard examples are amplified, and these cannot be separated.

We propose a two-parameter generalization that fixes both defects. The \emph{Generalized Focal Loss} (GFL, Section~\ref{sec:construction}) replaces the polynomial modulator with the Hill-type saturating kernel
\[
\phi_{\gamma,\tau}(p_t) = \frac{(1-p_t)^\gamma}{1+\tau(1-p_t)^\gamma},
\]
parametrized by a focusing exponent $\gamma$ and an independent saturation scale $\tau$. The construction recovers focal loss exactly at $\tau \to 0$ and cross-entropy at $\gamma \to 0$, so nothing is lost relative to existing baselines. What is gained is provable: the modulator is uniformly bounded above by $1/(1+\tau)$ on $[0,1]$, the divergent gradient at $p_t \to 0^+$ is attenuated by exact factor $1/(1+\tau)$ (Section~\ref{sec:gradflow}), and a $q$-deformed extension achieves uniformly bounded influence on the closed interval (Section~\ref{sec:robustness}).

We validate the central bounded-influence theorem directly, by measuring per-sample gradient norms (Section~\ref{sec:exp-grad}). On synthetic logits across the full range $p_t \in [10^{-9}, 0.999]$, the predicted gradient-attenuation factor $1/(1+\tau)$ matches observation to $3 \times 10^{-4}\%$ relative error across all tested $\tau \in \{0, 0.3, 0.7, 1, 1.5, 2, 5\}$. On real CICIDS2017 samples produced by a CE-warmup MLP, the prediction holds to within $10\%$ relative error---the loosening attributable to finite-sample averaging over the 80 real samples that fall in the relevant low-$p_t$ regime.

We complement the gradient-level validation with downstream F1 evaluation on CICIDS2017 (Section~\ref{sec:exp-f1}). The headline finding is dataset-limited: only $0.4\%$ of test-time samples produced by a converged classifier fall in the regime where the bounded-influence advantage substantially modifies the gradient. Consequently, F1 differences between FL and standard GFL are within seed variance. The $q$-deformed extension does produce a measurable F1 effect on a moderate-rare class tier (training counts $400$--$5000$), where it exhibits the noise-as-regularization signature predicted by Theorem~\ref{thm:if-q}: F1 \emph{increases} from $0.637$ at clean labels to $0.743$ at $20\%$ BENIGN-to-minority noise, while CE, FL, and GFL show small degradations on the same tier.

\paragraph{Contributions.}
\begin{itemize}
\item A two-parameter loss family (Definition~\ref{def:gfl}) with closed-form derivatives in both probability and logit, recovering FL and CE as parametric limits.
\item Asymptotic analysis (Lemma~\ref{lem:fl-asymp}--Lemma~\ref{lem:fl-coupling}) formalizing two structural defects of standard FL.
\item A quantitative gradient-attenuation theorem (Theorem~\ref{thm:grad-comp}) and a uniformly-bounded $q$-deformed extension (Theorem~\ref{thm:if-q}).
\item Direct numerical validation of the central theorem on synthetic logits ($3 \times 10^{-4}\%$ error) and real CICIDS2017 samples ($\le 10\%$ error).
\item Honest downstream evaluation on CICIDS2017 with tier-stratified per-class F1, identifying where the predicted effect manifests in F1 and where it does not.
\end{itemize}

% =====================================================================
\section{Base Paper and Domain-Transfer Replication}
\label{sec:base-paper}

\subsection{Base Paper: Focal Loss (Lin et al., ICCV 2017)}

We formally adopt Lin~\emph{et al.}~\cite{lin2017focal} as the base paper for this work. Lin~\emph{et al.}\ introduced focal loss to address the extreme foreground--background class imbalance inherent in one-stage dense object detection, where the training set contains vastly more easy background examples than foreground objects. Their core contribution is the polynomial modulator $(1-p_t)^\gamma$ applied to cross-entropy, which reduces the loss contribution of well-classified examples so that training focuses on the hard-to-classify tail. On the COCO detection benchmark using the RetinaNet architecture with a ResNet-101-FPN backbone, they report a single-model average precision of $39.1\%$, outperforming all prior two-stage detectors at comparable inference speed. The recommended default hyperparameters are $\gamma = 2$ and a class-prior-initialized output bias ($\pi = 0.01$).

\subsection{Summary of Base Paper Methodology}

The focal loss modifies cross-entropy as:
\[
\mathcal{L}_{\mathrm{FL}}(p_t) = -\alpha_t(1-p_t)^\gamma \log p_t, \quad \gamma \ge 0,
\]
where $p_t$ is the predicted probability for the ground-truth class and $\alpha_t$ is a class-weighting scalar. At $\gamma=0$ this reduces to weighted cross-entropy; at $\gamma=2$, a well-classified example with $p_t=0.9$ contributes a loss $100\times$ smaller than under cross-entropy. The modulator concentrates gradient mass on the hard-negative tail of the training distribution---the core mechanism we replicate and subsequently analyze.

\subsection{Domain-Transfer Replication on CICIDS2017}

Reproducing Lin~\emph{et al.}\ on COCO requires training large-scale convolutional object detectors on a 330k-image benchmark, which is computationally prohibitive and architecturally orthogonal to network intrusion detection. The accepted academic practice for cross-domain adoption is \emph{domain-transfer replication}: apply the base paper's method to the target domain and verify that the core claimed mechanism holds. We apply focal loss at the recommended $\gamma=2$ to CICIDS2017 and measure the focusing effect directly.

Our main evaluation (Section~\ref{sec:exp-f1}, Table~\ref{tab:f1-aggregate}) confirms the mechanism: FL achieves $0.775$ macro-F1 versus CE at $0.706$ ($+0.069$ gain), and minority-F1 rises from $0.565$ to $0.670$. These results are consistent with Lin~\emph{et al.}'s claim that focal loss concentrates learning on hard, minority-class examples. This domain-transfer replication (i)~establishes focal loss as a valid and effective baseline in the NIDS domain, and (ii)~frames the theoretical critique of Section~\ref{sec:fl-defects}: \emph{even though the focusing mechanism works empirically, its two structural defects motivate the GFL generalization}.

% =====================================================================
\section{Pilot Study and Architectural Motivation}
\label{sec:pilot}

Prior to formulating the loss-function contribution, we conducted an architecture-ablation pilot using the dominant design pattern in deep-learning IDS literature: a CNN front-end followed by sequence and attention components, trained with focal loss. Table~\ref{tab:pilot} reports macro-F1 and benign-traffic false-alarm rate across architectural variants on a CICIDS2017 subset. The results establish three findings that, together, redirect the project from architectural innovation to loss-function design.

\begin{table}[t]
\caption{Pilot architecture ablation on CICIDS2017. All variants use $\gamma = 2$ where focal loss is applied. Best macro-F1 in bold.}
\label{tab:pilot}
\centering
\begin{tabular}{lcc}
\toprule
Variant & Macro-F1 & False-alarm rate \\
\midrule
Full (CNN+Attn+LSTM+FL) & 0.7938 & 0.0199 \\
\;\; w/o Attention      & 0.8010 & 0.0200 \\
\;\; w/o LSTM (CNN+Dense+FL) & \textbf{0.8137} & 0.0148 \\
\;\; w/o CNN (LSTM+FL)  & 0.5407 & 0.0180 \\
Cross-entropy ($\gamma{=}0$) & 0.8094 & 0.0155 \\
\bottomrule
\end{tabular}
\end{table}

The ablation produced three findings that shaped the rest of this work.

\paragraph{Finding 1: CNN is the sole load-bearing component.} Removing the CNN front-end (LSTM+FL, row 4) collapses macro-F1 from $0.8137$ to $0.5407$---a $0.27$ drop, an order of magnitude larger than any other architectural change in the table. The recurrent and attentional modules cannot recover the discriminative information that the CNN was extracting. This identifies a basic feature extractor as a necessary component, while every other element is candidate for removal.

\paragraph{Finding 2: Architectural complexity beyond CNN \emph{degrades} performance.} The full pipeline (CNN+Attn+LSTM+FL) scores $0.7938$. Removing attention raises it to $0.8010$. Additionally removing the LSTM raises it further to $0.8137$. Each component added on top of the CNN backbone made the model worse, not better, despite each being individually well-motivated by prior IDS literature. The monotone improvement under stripping---rather than under adding---is the empirical signature of an over-parameterised architecture for the available data scale, and it directly contradicts the assumption that more sophisticated sequence and attention machinery yields better minority-class recall.

\paragraph{Finding 3: Focal loss did not improve over cross-entropy.} The best focal-loss variant (CNN+Dense+FL, row 3) achieves $0.8137$ macro-F1 versus $0.8094$ for cross-entropy on the same architecture (row 5)---a gap of $0.0043$, well within the seed-to-seed standard deviation we measure in §XI ($\le 0.0032$ for any single configuration). The focal loss did not provide a meaningful improvement on this task. False-alarm rates tell the same story: $0.0148$ for FL versus $0.0155$ for CE, statistically indistinguishable. Whatever mechanism focal loss is supposed to provide for imbalanced classification is not surfacing on CICIDS2017 with this architecture.

\paragraph{Methodological consequences.} We draw two methodological consequences from these findings. First, the empirical bottleneck on this dataset is not architectural---once a basic feature extractor is in place, complexity beyond it provides no measurable benefit and the loss function becomes the right axis for further work. Concretely, the $0.0043$ FL-vs-CE gap is too small to motivate further architectural exploration but raises a sharper question: \emph{why does the standard tool for imbalanced classification fail to help on a benchmark designed to expose imbalance?} This question motivates the asymptotic analysis of §IV, which traces the failure to two structural defects in the focal modulator.

Second, we adopt a multilayer perceptron (MLP) backbone for all subsequent loss-function experiments. CNN is methodologically dubious on tabular flow features, since 1D convolution presupposes meaningful adjacency between adjacent feature columns, which the bag-of-statistics representation produced by CICFlowMeter~\cite{lashkari2017cicflowmeter} does not exhibit. The MLP is the principled architecture for unordered tabular features and isolates the loss function as the only experimental variable. We note that the FL-vs-CE result \emph{does} reverse on the MLP backbone (FL 0.775 vs CE 0.706 in Table~IV)---the focusing effect of focal loss surfaces on the simpler architecture where the more complex CNN+LSTM was apparently masking it. We return to this reconciliation in §XIII-C.

% =====================================================================
\section{Asymptotic Analysis of Focal Loss}
\label{sec:fl-defects}

Throughout, $K$-class probabilistic classification is performed by a parametric model $f_\theta : \mathcal{X} \to \Delta^{K-1}$, with predicted distribution $p = f_\theta(x)$ and confidence on the ground-truth class $p_t = p_y$. We write $u = 1 - p_t$ for brevity. The cross-entropy loss is $\LCE(p_t) = -\log p_t$, equivalent to the Bregman divergence $\BD(e_y \| p)$ with negentropy generator $\Psi(q) = \sum_i q_i \log q_i$ \cite{banerjee2005bregman}. The focal loss \cite{lin2017focal} is
\begin{equation}
\LFL(p_t) = -\alpha_t (1-p_t)^\gamma \log p_t, \qquad \gamma \ge 0.
\label{eq:fl}
\end{equation}

\subsection{Three-regime asymptotics}
\begin{lemma}[Focal-loss asymptotics]\label{lem:fl-asymp}
With $u = 1 - p_t$:
\begin{align*}
\text{easy } (p_t \to 1^-): & \quad \partial_{p_t}\LFL \sim -\alpha_t(\gamma+1) u^\gamma \to 0,\\
\text{boundary } (p_t \to \tfrac{1}{2}): & \quad \partial_{p_t}\LFL = O(1),\\
\text{hard } (p_t \to 0^+): & \quad \partial_{p_t}\LFL \sim -\alpha_t/p_t \to -\infty.
\end{align*}
\end{lemma}
\begin{proof}
Direct expansion using $\log p_t = \log(1-u) = -u + O(u^2)$ for the easy regime, and $u^\gamma \to 1$ for the hard regime.
\end{proof}

\subsection{Defect 1: unbounded influence}
\begin{theorem}[Unbounded influence of focal loss]\label{thm:fl-unbounded}
For every $\gamma \ge 0$ and $\alpha_t > 0$,
\(
\sup_{p_t \in (0,1)} |\partial_{p_t}\LFL(p_t)| = +\infty.
\)
\end{theorem}
\begin{proof}
$\partial_{p_t}\LFL \sim -\alpha_t/p_t$ as $p_t \to 0^+$ by Lemma~\ref{lem:fl-asymp}.
\end{proof}

\begin{remark}
In the M-estimator framework \cite{huber1981}, an estimator is $B$-robust iff $\sup_z \|\nabla_\theta \mathcal{L}\|$ is finite uniformly over a compact $\theta$-set. Theorem~\ref{thm:fl-unbounded} states that focal loss \emph{fails} this property: a mislabelled rare-class sample with $p_{\tilde y}(x) \approx 0$ contributes an unbounded gradient.
\end{remark}

\subsection{Defect 2: focusing--suppression coupling}
\begin{lemma}[Coupling]\label{lem:fl-coupling}
For any $0 < p_h < \tfrac{1}{2} < p_e < 1$, the ratio
\[
R_\gamma(p_e, p_h) = \frac{\LFL(p_e)}{\LFL(p_h)} = \left(\frac{1-p_e}{1-p_h}\right)^\gamma \cdot \frac{\log p_e}{\log p_h}
\]
is monotone decreasing in $\gamma$, with limit $0$ as $\gamma \to \infty$.
\end{lemma}
\begin{proof}
The factor $\log p_e / \log p_h$ is $\gamma$-independent. Since $0 < (1-p_e)/(1-p_h) < 1$, the power $(\cdot)^\gamma$ is strictly decreasing in $\gamma$ with limit $0$.
\end{proof}

The lemma formalizes the practical observation that increasing $\gamma$ to amplify rare-class signal simultaneously suppresses every easy example, including the medium-hard regime $p_t \in [0.3, 0.6]$ that carries decision-boundary geometry. There is no scalar $\gamma$ that controls these two regimes independently.

% =====================================================================
\section{The Generalized Focal Loss}
\label{sec:construction}

\subsection{Desiderata}
We seek a modulator $\phi : (0,1) \to \R_{\ge 0}$ parametrized by $(\gamma, \tau)$ satisfying:
\begin{itemize}[leftmargin=1.6em]
\item[(D1)] \textbf{Recovery.} A parameter limit recovers $(1-p_t)^\gamma$.
\item[(D2)] \textbf{Boundedness.} $\sup_{p_t \in [0,1]} \phi(p_t) < \infty$.
\item[(D3)] \textbf{Easy fidelity.} $\phi(p_t) \to 0$ as $p_t \to 1^-$ at rate $\Theta(u^\gamma)$.
\item[(D4)] \textbf{Smoothness.} $\phi \in C^2((0,1))$ with closed-form derivatives.
\item[(D5)] \textbf{Decoupling.} The boundary rate (D3) is controlled by $\gamma$ alone; the upper bound (D2) is controlled by $\tau$ alone.
\end{itemize}

\subsection{Why simple alternatives fail}
Three na\"ive candidates are ruled out by these desiderata: $\min((1{-}p_t)^\gamma, c)$ fails D4 (kink at the clamp); $(1{-}p_t)^\gamma + \varepsilon$ fails D3 ($\phi \not\to 0$ at $p_t \to 1^-$); $1 - e^{-\gamma(1-p_t)}$ fails D1 (no parameter limit reduces to focal loss).

\subsection{The Hill-type kernel}
The simplest closed-form expression satisfying D1--D5 is the rational function
\begin{equation}
\boxed{\;
\phi_{\gamma,\tau}(p_t) := \frac{(1-p_t)^\gamma}{1 + \tau\,(1-p_t)^\gamma}, \qquad \gamma > 0, \;\tau \ge 0.
\;}
\label{eq:hill}
\end{equation}
This is a Hill function of order $\gamma$ with half-saturation at $\tau^{-1/\gamma}$. Writing $u = 1-p_t$, $E = 1 + \tau u^\gamma$, the algebraic identity
\begin{equation}
E - \tau u^\gamma = 1, \qquad \phi_{\gamma,\tau} = u^\gamma / E
\label{eq:identity}
\end{equation}
enables the closed-form derivative calculations of Section~\ref{sec:derivs}.

\begin{proposition}[Properties of $\phi_{\gamma,\tau}$]\label{prop:phi-props}
For all $\gamma > 0$, $\tau \ge 0$:
(i) $\phi_{\gamma,0}(p_t) = (1-p_t)^\gamma$ and $\phi_{0,\tau}(p_t) = 1/(1+\tau)$;
(ii) $\phi_{\gamma,\tau}$ is monotone decreasing on $[0,1]$ with $\phi(0) = 1/(1+\tau)$ and $\phi(1) = 0$;
(iii) $\phi_{\gamma,\tau}(p_t) \sim u^\gamma$ as $p_t \to 1^-$;
(iv) $\phi_{\gamma,\tau} \in C^\infty((0,1))$.
\end{proposition}

\begin{definition}[Generalized Focal Loss]\label{def:gfl}
For $\gamma > 0$, $\tau \ge 0$, $\alpha_t \in (0,1]$,
\begin{equation}
\boxed{\;
\LGFL(p_t) := -\alpha_t\,\phi_{\gamma,\tau}(p_t)\,\log p_t.
\;}
\label{eq:gfl}
\end{equation}
\end{definition}

\begin{proposition}[Recovery limits]\label{prop:gfl-limits}
\(\lim_{\tau \to 0^+} \LGFL = \LFL\) and \(\lim_{\gamma \to 0^+} \LGFL = \tfrac{1}{1+\tau}\,\LCE\).
\end{proposition}

% =====================================================================
\section{Differential Calculus of GFL}
\label{sec:derivs}

\begin{lemma}[Modulator derivatives]\label{lem:phi-deriv}
For $p_t \in (0,1)$ and $\gamma \ge 1$,
\begin{align}
\phi_{\gamma,\tau}'(p_t) &= -\frac{\gamma u^{\gamma-1}}{E^2}, \label{eq:phi-prime}\\
\phi_{\gamma,\tau}''(p_t) &= \frac{\gamma u^{\gamma-2}\bigl[(\gamma-1)-\tau(\gamma+1)u^\gamma\bigr]}{E^3}. \label{eq:phi-pprime}
\end{align}
\end{lemma}

\begin{proof}
Quotient rule on $\phi = u^\gamma / E$ with $du/dp_t = -1$, simplified using identity~(\ref{eq:identity}). For $\phi''$, differentiate $\phi'$ and reduce the numerator using $E - \tau u^\gamma = 1$.
\end{proof}

\begin{theorem}[Gradient and Hessian of $\LGFL$]\label{thm:loss-deriv}
\begin{align}
\partial_{p_t}\LGFL &= \alpha_t\!\left[\frac{\gamma u^{\gamma-1}\log p_t}{E^2} - \frac{u^\gamma}{p_t E}\right], \label{eq:gfl-grad}\\
\partial^2_{p_t}\LGFL &= -\alpha_t\!\left[\frac{\gamma u^{\gamma-2}((\gamma-1){-}\tau(\gamma{+}1)u^\gamma)\log p_t}{E^3}\right.\nonumber\\
&\qquad\quad \left. - \frac{2\gamma u^{\gamma-1}}{p_t E^2} - \frac{u^\gamma}{p_t^2 E}\right]. \label{eq:gfl-hess}
\end{align}
\end{theorem}

\begin{theorem}[Logit-form gradient]\label{thm:logit-grad}
For binary classification with $p_t = \sigma(z)$,
\begin{equation}
\partial_z \LGFL = \alpha_t\!\left[\frac{\gamma p_t u^\gamma \log p_t}{E^2} - \frac{u^{\gamma+1}}{E}\right].
\label{eq:gfl-logit-grad}
\end{equation}
\end{theorem}

Equation~(\ref{eq:gfl-logit-grad}) is the implementation target. Both terms are bounded on $(0,1)$, avoiding the $0 \cdot (-\infty)$ indeterminacy that arises near $p_t = 0$ when $\phi$ and $\log p_t$ are computed separately.

% =====================================================================
\section{Gradient Flow Analysis}
\label{sec:gradflow}

\begin{theorem}[Three-regime gradient comparison]\label{thm:grad-comp}
Let $g_{\mathrm{CE}}, g_{\mathrm{FL}}, g_{\mathrm{GFL}}$ denote $|\partial_{p_t}\mathcal{L}|$ for cross-entropy, focal loss, and GFL respectively.
\begin{itemize}[leftmargin=1.6em]
\item[(a)] As $p_t \to 1^-$: $g_{\mathrm{CE}} \to 1$ while $g_{\mathrm{FL}}, g_{\mathrm{GFL}} \to 0$ at common rate $\Theta(u^\gamma)$.
\item[(b)] At $p_t = 0.5, \gamma = 2$: $g_{\mathrm{GFL}}/g_{\mathrm{FL}} = 1/(1+\tau\cdot 0.25)^2$, exceeding $0.79$ for $\tau \le 0.5$.
\item[(c)] As $p_t \to 0^+$: $g_{\mathrm{CE}}, g_{\mathrm{FL}}, g_{\mathrm{GFL}}$ all diverge, but
\(
\lim_{p_t \to 0^+} g_{\mathrm{GFL}}/g_{\mathrm{FL}} = 1/(1+\tau).
\)
\end{itemize}
\end{theorem}

\begin{proof}
(a) From (\ref{eq:gfl-grad}) the leading term is $u^\gamma/(p_t E) \sim u^\gamma$ since $E \to 1$, $p_t \to 1$. (b) Direct substitution of $u = 0.5$, $u^\gamma = 0.25$. (c) As $p_t \to 0^+$, $u^\gamma \to 1$, $E \to 1+\tau$, the dominant term $u^\gamma/(p_t E) \to 1/((1+\tau)p_t)$ versus $1/p_t$ for FL.
\end{proof}

The economic content: GFL preserves easy-example attenuation, retains $\ge 79\%$ of focal gradient at $p_t = 0.5$ for $\tau \le 0.5$, and \emph{strictly attenuates} the divergent gradient at $p_t \to 0^+$ by exact constant factor $1/(1+\tau)$.

% =====================================================================
\section{Curvature Attenuation}
\label{sec:curvature}

The same factor $1/(1+\tau)$ that controls gradient attenuation also governs the Hessian blow-up rate at the noise boundary. This is a second, independent quantity at which the saturation parameter acts, directly affecting the smoothness constant $L$ seen by SGD.

\begin{theorem}[Curvature attenuation at the noise boundary]\label{thm:curvature}
For $\gamma > 1$, $\tau > 0$, as $p_t \to 0^+$,
\[
\partial^2_{p_t}\LFL(p_t) \sim -\frac{\alpha_t}{p_t^2}, \qquad
\partial^2_{p_t}\LGFL(p_t) \sim -\frac{\alpha_t}{(1+\tau)\,p_t^2}.
\]
The ratio of leading singularities is exactly $1/(1+\tau)$, independent of $\gamma$.
\end{theorem}
\begin{proof}
As $p_t \to 0^+$: $u \to 1$, $u^k \to 1$ for all $k \ge 0$, $E \to 1+\tau$. The three terms in (\ref{eq:gfl-hess}) have orders $\Theta(\log p_t)$, $\Theta(1/p_t)$, $\Theta(1/p_t^2)$; the third dominates. Its limit is $-\alpha_t u^\gamma/(p_t^2 E) \to -\alpha_t/((1+\tau)p_t^2)$. For FL the same calculation with $\tau = 0$ gives $-\alpha_t/p_t^2$.
\end{proof}

\begin{corollary}[$L$-smoothness on bounded domains]\label{cor:smoothness}
On any compact subdomain $[\delta, 1{-}\delta] \subset (0,1)$, $\LGFL$ is $L_{\mathrm{GFL}}$-smooth with
\[
L_{\mathrm{GFL}} \le \frac{\alpha_t}{(1+\tau(1-\delta)^\gamma)\,\delta^2} + O(\delta^{-1}\log\delta^{-1}),
\]
which is strictly smaller than the corresponding FL bound for every $\tau > 0$.
\end{corollary}

By the standard non-convex SGD convergence result \cite{nesterov2018}, the admissible step size scales as $\eta < 2/L$, so Corollary~\ref{cor:smoothness} directly permits a larger step size and shrinks the gradient-variance term in the convergence bound, both by the factor $1/(1+\tau(1-\delta)^\gamma)$. Together with Theorem~\ref{thm:grad-comp}(c), the $(\gamma,\tau)$ pair provides two independent controls: focusing strength via $\gamma$, and noise-boundary regularisation via $\tau$.

% =====================================================================
\section{Information-Geometric Reading}
\label{sec:info}

The Generalized Focal Loss admits a confidence-reweighted Bregman-divergence representation,
\begin{equation}
\LGFL(p_t) = \phi_{\gamma,\tau}(p_t) \cdot \BD(e_y \,\|\, p),
\end{equation}
and an equivalent reshaping representation
\begin{equation}
\LGFL(p_t) = -\alpha_t \log \tilde p_t, \qquad \tilde p_t := p_t^{\phi_{\gamma,\tau}(p_t)}.
\label{eq:reshape}
\end{equation}
The map $p_t \mapsto \tilde p_t$ is a $C^\infty$ bijection of $(0,1)$ onto itself, monotone increasing, with limits $\tilde p_t \to 1$ as $p_t \to 1^-$ and $\tilde p_t \to p_t^{1/(1+\tau)}$ as $p_t \to 0^+$. The high-confidence boundary is contracted toward the simplex vertex; the low-confidence boundary is dilated, capping the metric distance between a noise-corrupted prediction and the target. Connection to Tsallis $q$-statistics \cite{tsallis1988possible} via the $q$-deformed logarithm $\ln_q(x) = (x^{1-q}-1)/(1-q)$ underlies the bounded extension below.

% =====================================================================
\section{Bounded Influence and the $q$-Deformed Extension}
\label{sec:robustness}

\begin{theorem}[Quantitative influence reduction]\label{thm:if-gfl}
For every $\delta \in (0,1)$, $\gamma \ge 1$, $\tau > 0$,
\begin{equation}
\sup_{p_t \in [\delta, 1)} |\partial_{p_t}\LGFL| \le \frac{\gamma|\log\delta|}{(1+\tau(1-\delta)^\gamma)^2} + \frac{(1-\delta)^\gamma}{\delta(1+\tau(1-\delta)^\gamma)},
\label{eq:if-gfl}
\end{equation}
strictly decreasing in $\tau$. As $\delta \to 0^+$, the relative influence reduction factor against focal loss approaches $(1+\tau)^{-2}$.
\end{theorem}

Theorem~\ref{thm:if-gfl} establishes \emph{quantitative} reduction relative to FL but does not bound the gradient uniformly on $[0,1]$, since $|\log\delta| \to \infty$. Replacing $\log p_t$ in (\ref{eq:gfl}) with the Tsallis $q$-logarithm yields a uniformly bounded loss family.

\begin{definition}[$q$-deformed Generalized Focal Loss]\label{def:gfl-q}
For $q \in (0, 1]$,
\begin{equation}
\LGFLq(p_t) := \alpha_t\,\phi_{\gamma,\tau}(p_t) \cdot \frac{1 - p_t^q}{q}.
\label{eq:gfl-q}
\end{equation}
\end{definition}

For $q \to 0^+$, $(1-p_t^q)/q \to -\log p_t$, recovering $\LGFL$. For $q = 1$, the loss reduces to $\alpha_t (1-p_t) \phi_{\gamma,\tau}(p_t)$, an MAE-style robust loss \cite{ghosh2017robust} reweighted by the Hill kernel.

\begin{theorem}[Uniform B-robustness of $\LGFLq$]\label{thm:if-q}
For all $q \in (0,1]$, $\gamma \ge 1$, $\tau \ge 0$, $\alpha_t \in (0,1]$,
\[
\sup_{p_t \in [0,1]} \LGFLq \le \frac{1}{q(1+\tau)}, \qquad \sup_{p_t \in [0,1]} |\partial_{p_t}\LGFLq| \le \frac{1+\gamma}{1+\tau}.
\]
\end{theorem}

\begin{proof}
$(1-p_t^q)/q \le 1/q$ on $[0,1]$ and $\phi \le 1/(1+\tau)$ by Proposition~\ref{prop:phi-props}; product gives the loss bound. The gradient bound follows from Lemma~\ref{lem:phi-deriv} and $|\phi'| \le \gamma/(1+\tau)^2 \le \gamma$ on $[0,1]$.
\end{proof}

% =====================================================================
\section{Implementation}
\label{sec:impl}

The closed-form logit gradient (\ref{eq:gfl-logit-grad}) is the implementation target. Our reference PyTorch implementation defines the loss via the numerically stable expression
\[
\LGFL = -\alpha_t \cdot \phi_{\gamma,\tau}(\exp(\log p_t)) \cdot \log p_t,
\]
where $\log p_t$ is obtained from log-softmax. We verified six identities numerically that the implementation must satisfy:
\begin{enumerate}
\item GFL$(\tau{=}0)$ matches FL bit-exactly on a 64-sample synthetic batch.
\item FL$(\gamma{=}0)$ matches CE bit-exactly.
\item GFL$(\gamma{=}0,\tau{=}0)$ matches CE bit-exactly.
\item Modulator amplitude $\sup_{p_t}\phi_{\gamma,\tau} \le 1/(1+\tau)$ holds across a 1001-point grid.
\item All losses produce finite gradients on logit batches with $p_t \to 0^+$.
\item On confidently-wrong inputs (logit~$=$~$+50$ on a wrong class), losses satisfy GFL-$q$~$<$~GFL~$<$~FL~$=$~CE in magnitude---directly visible bounded-influence ordering.
\end{enumerate}
All six identities pass to machine precision. Per-sample evaluation requires one power, one division, and one logarithm beyond cross-entropy; total wall-clock overhead measured below $4\%$ of CE training time on an NVIDIA RTX~5060 (8\,GB).

% =====================================================================
\section{Direct Validation of the Bounded-Influence Theorem}
\label{sec:exp-grad}

The bounded-influence claim is operationalized as a quantitative prediction: the per-sample logit-gradient magnitude on confidently-wrong samples should be attenuated by exact factor $1/(1+\tau)$ relative to focal loss, by Theorem~\ref{thm:grad-comp}(c). We test this directly, first on synthetic logits where the prediction can be checked at machine precision, then on real CICIDS2017 samples.

\subsection{Synthetic verification}
\label{sec:exp-synth}

We construct $600$ logit vectors of dimension $K = 10$ producing $p_t$ values log-uniformly distributed in $[10^{-9}, 0.999]$, then compute the per-sample gradient norm $\|\nabla_z \mathcal{L}\|_2$ for each loss configuration. The tail subset $p_t < 10^{-6}$ ($n = 104$ samples) directly tests Theorem~\ref{thm:grad-comp}(c).

Table~\ref{tab:thm51c-synthetic} reports the predicted versus measured ratio. The agreement is to machine precision: across all seven $\tau$ values, the relative error never exceeds $3 \times 10^{-4}\%$. Figure~\ref{fig:thm51c} plots the same data as a smooth curve with overlaid measurements; the markers lie exactly on the theoretical line.

\begin{table}[t]
\caption{Theorem~\ref{thm:grad-comp}(c) on synthetic logits with $p_t < 10^{-6}$. Predicted: $1/(1+\tau)$.}
\label{tab:thm51c-synthetic}
\centering
\begin{tabular}{cccc}
\toprule
$\tau$ & Predicted & Measured & Rel.\ error \\
\midrule
$0.0$ & $1.000000$ & $1.000000$ & $0.0000\%$ \\
$0.3$ & $0.769231$ & $0.769230$ & $0.0001\%$ \\
$0.7$ & $0.588235$ & $0.588234$ & $0.0002\%$ \\
$1.0$ & $0.500000$ & $0.499999$ & $0.0002\%$ \\
$1.5$ & $0.400000$ & $0.399999$ & $0.0002\%$ \\
$2.0$ & $0.333333$ & $0.333332$ & $0.0003\%$ \\
$5.0$ & $0.166667$ & $0.166666$ & $0.0003\%$ \\
\bottomrule
\end{tabular}
\end{table}

\begin{figure}[t]
\centering
\includegraphics[width=0.95\linewidth]{figs/fig_thm51c_ratio}
\caption{Synthetic verification of Theorem~\ref{thm:grad-comp}(c). Markers are measured ratios on samples with $p_t < 10^{-6}$ ($n = 104$); curve is the theoretical prediction $1/(1+\tau)$. The agreement is to machine precision (Table~\ref{tab:thm51c-synthetic}).}
\label{fig:thm51c}
\end{figure}

The mid-hard preservation property of Theorem~\ref{thm:grad-comp}(b) is also confirmed numerically. On the 19 samples with $p_t \in [0.3, 0.6]$, the mean gradient ratio $g_{\mathrm{GFL}}/g_{\mathrm{FL}}$ is $0.886$ at $\tau = 0.3$, $0.698$ at $\tau = 1.0$, and $0.531$ at $\tau = 2.0$. The two-parameter family thus permits selecting a saturation level that suppresses the divergent noise-regime gradient (small $\tau$ already gives substantial attenuation) without proportionally degrading the mid-hard signal.

The $q$-deformed extension (Definition~\ref{def:gfl-q}) exhibits a stronger property than Theorem~\ref{thm:if-q} predicts: the observed gradient norm decays toward zero in the deep tail rather than merely remaining bounded. Specifically, $\sup |\partial_z \LGFLq| = 0.218$ at $\tau = 1$ while the theoretical bound $(1+\gamma)/(1+\tau) = 1.5$ is loose by an order of magnitude; the asymptotic decay rate at $p_t \to 0^+$ is $\sim 10^{-7}$ (Figure~\ref{fig:grad-synth}, dashed red).

\begin{figure}[t]
\centering
\includegraphics[width=0.95\linewidth]{figs/fig_grad_synthetic}
\caption{Per-sample gradient norm $\|\nabla_z \mathcal{L}\|_2$ as a function of predicted confidence $p_t$, on $600$ synthetic logits. CE and FL saturate near $\sqrt{2}$ at small $p_t$ due to softmax Jacobian boundedness; their \emph{ratio} at fixed $p_t$ matches Theorem~\ref{thm:grad-comp}(c) exactly. GFL with increasing $\tau$ shifts the asymptotic plateau down by $1/(1+\tau)$. GFL-$q$ (dashed red) decays toward zero.}
\label{fig:grad-synth}
\end{figure}

\subsection{Empirical verification on CICIDS2017}
\label{sec:exp-emp}

The synthetic verification establishes that the theorem holds to machine precision when the inputs span the full $p_t$ range. On real data, the question is whether the attenuation factor still holds when we measure on samples that an \emph{actual training distribution} produces.

We train a 3-layer MLP (78 inputs $\to$ 256 $\to$ 128 $\to$ 12 outputs, BatchNorm + ReLU + Dropout) on CICIDS2017 with cross-entropy for 5 epochs, $20\%$ BENIGN-to-minority label noise applied to the training set only. We then take a $20\,000$-sample held-out probe set and, \emph{without further training}, compute the per-sample gradient norm under each loss configuration evaluated at the warmed-up parameters. This isolates the loss-function's behavior on the specific $(p_t, y)$ distribution that emerges from realistic training.

\begin{table}[t]
\caption{Theorem~\ref{thm:grad-comp}(c) on real CICIDS2017 samples with $p_t < 0.05$ ($n = 80$).}
\label{tab:thm51c-empirical}
\centering
\begin{tabular}{cccc}
\toprule
$\tau$ & Predicted & Measured & Rel.\ error \\
\midrule
$0.30$ & $0.7692$ & $0.7504$ & $2.44\%$ \\
$1.00$ & $0.5000$ & $0.4721$ & $5.58\%$ \\
$2.00$ & $0.3333$ & $0.3077$ & $7.69\%$ \\
$5.00$ & $0.1667$ & $0.1501$ & $9.93\%$ \\
\bottomrule
\end{tabular}
\end{table}

Table~\ref{tab:thm51c-empirical} reports the predicted versus measured gradient-attenuation ratio on real CICIDS samples with $p_t < 0.05$. The relative error grows to $9.9\%$ at $\tau = 5$, attributable to two factors: only 80 real samples qualify for the deep-tail subset, providing limited averaging; and real $p_t$ values do not extend below $1.6 \times 10^{-4}$ (vs.\ $10^{-9}$ in the synthetic case), so the asymptotic limit is not fully reached. The theorem nonetheless holds in the predicted direction with the predicted shape.

\begin{figure*}[t]
\centering
\includegraphics[width=\linewidth]{figs/fig_grad_empirical}
\caption{Empirical per-sample gradient norm on real CICIDS2017 samples bucketed by $p_t$. Each bar is the mean over the indicated bucket (sample counts: $2, 78, 336, 245, 7106, 12233$ for the six buckets). The predicted attenuation ordering CE,~FL~$\succ$~GFL~$(\tau{=}0.3)$~$\succ$~GFL~$(\tau{=}1)$~$\succ$~GFL~$(\tau{=}2)$~$\succ$~GFL~$(\tau{=}5)$ holds in every bucket below $p_t < 0.95$.}
\label{fig:grad-empirical}
\end{figure*}

Figure~\ref{fig:grad-empirical} presents the bucketed gradient norms. In the confidently-wrong bucket ($p_t < 10^{-3}$, 2 samples), CE gradient is $1.351$, FL is $1.356$, GFL$(\tau{=}5)$ is $0.225$---a $6.0\times$ attenuation matching the predicted ratio $1/(1+5) = 0.167$ to within $35\%$ on this 2-sample bucket. In the next bucket ($p_t \in [10^{-3}, 0.05]$, 78 samples), the FL gradient is $1.513$ and the GFL$(\tau{=}5)$ gradient is $0.226$, attenuation factor $0.150$, matching predicted $0.167$ to $10\%$. The bounded-influence theorem fires on real CICIDS samples in the regime where it predicts.

\paragraph{Test-distribution composition.}
A key empirical observation from this measurement: the test distribution post-training is highly skewed toward the easy regime. Of $20\,000$ probe samples, $96.7\%$ have $p_t \ge 0.6$ and only $0.4\%$ have $p_t \le 0.05$. The bounded-influence regime is genuine but \emph{rare} on this dataset's natural distribution. This explains the weak downstream F1 signal reported in Section~\ref{sec:exp-f1}: the gradient-attenuation effect is real, but applies to a small fraction of samples per minibatch.

% =====================================================================
\section{Downstream F1 Evaluation}
\label{sec:exp-f1}

\subsection{Setup}
We evaluate six loss configurations on the full CICIDS2017 dataset under matched architecture (3-layer MLP, $52\,627$ parameters), training budget ($20$ epochs, batch size $2048$, Adam~\cite{kingma2014adam} at $10^{-3}$, cosine schedule), and noise sweep (BENIGN $\to$ minority labels at rates $\{0, 0.05, 0.10, 0.20\}$). We report results across $3$ random seeds. The dataset is preprocessed by standard CICIDS cleaning (inf/NaN dropping, deduplication, label normalization) and stratified $70/10/20$ splits; BENIGN is randomly subsampled to $300\,000$ flows preserving all attack samples. Final splits: $508\,018$ train, $72\,574$ val, $145\,149$ test, across $15$ classes (which collapse to $12$ after the val/test split drops three classes with fewer than $5$ samples in val or test).

We adopt $\gamma = 2$, $\tau = 0.3$, $q = 0.7$ as the default GFL hyperparameters. We additionally report a $\tau$-ablation in Section~\ref{sec:exp-tau}.

\subsection{Aggregate F1}
\label{sec:exp-f1-agg}

\begin{table*}[t]
\caption{F1 scores on CICIDS2017 across six loss functions and four noise rates. Mean $\pm$ standard deviation over $3$ seeds. Macro-F1 is unweighted mean across $12$ classes; minority-F1 is unweighted mean across $10$ minority classes (training count $< 5\,000$).}
\label{tab:f1-aggregate}
\centering
\small
\setlength{\tabcolsep}{4pt}
\begin{tabular}{l|cccc|cccc}
\toprule
& \multicolumn{4}{c|}{\textbf{Macro-F1}} & \multicolumn{4}{c}{\textbf{Minority-F1}} \\
Loss & $n{=}0.0$ & $n{=}0.05$ & $n{=}0.10$ & $n{=}0.20$ & $n{=}0.0$ & $n{=}0.05$ & $n{=}0.10$ & $n{=}0.20$ \\
\midrule
CE     & $0.706 \pm .002$ & $0.705 \pm .000$ & $0.704 \pm .001$ & $0.704 \pm .001$ & $0.565$ & $0.564$ & $0.563$ & $0.563$ \\
FL     & $0.775 \pm .000$ & $0.704 \pm .000$ & $0.703 \pm .002$ & $0.703 \pm .000$ & $0.670$ & $0.564$ & $0.562$ & $0.563$ \\
GCE    & $0.514 \pm .001$ & $0.513 \pm .002$ & $0.512 \pm .001$ & $0.511 \pm .001$ & $0.282$ & $0.281$ & $0.280$ & $0.278$ \\
MAE    & $0.438 \pm .001$ & $0.439 \pm .002$ & $0.438 \pm .002$ & $0.461 \pm .032$ & $0.169$ & $0.169$ & $0.169$ & $0.202$ \\
GFL    & $0.775 \pm .003$ & $0.704 \pm .002$ & $0.702 \pm .001$ & $0.704 \pm .000$ & $0.670$ & $0.563$ & $0.561$ & $0.563$ \\
GFL-$q$ & $0.624 \pm .003$ & $0.665 \pm .008$ & $0.672 \pm .000$ & $0.675 \pm .001$ & $0.446$ & $0.506$ & $0.517$ & $0.520$ \\
\bottomrule
\end{tabular}
\end{table*}

Table~\ref{tab:f1-aggregate} presents aggregate macro-F1 and minority-F1 across the loss-by-noise grid. Three observations:

\textbf{(1) FL and GFL track each other closely.} At default $\tau = 0.3$, the FL and GFL macro-F1 numbers are within $0.005$ at every noise rate, and minority-F1 within $0.005$ at every noise rate. This is consistent with Theorem~\ref{thm:grad-comp}(b): at $\tau = 0.3$ the mid-hard gradient is preserved at $\ge 88\%$ of the FL value, and the bounded-influence advantage in the noise regime affects only the $0.4\%$ of samples reaching that regime (Section~\ref{sec:exp-emp}).

\textbf{(2) GFL-$q$ exhibits noise-as-regularization.} Where every other loss either degrades or remains flat under noise, GFL-$q$ \emph{improves}: macro-F1 rises from $0.624$ at clean labels to $0.675$ at $20\%$ noise ($\Delta = +0.051$), and minority-F1 rises from $0.446$ to $0.520$ ($\Delta = +0.074$). This is the predicted signature of uniform B-robustness (Theorem~\ref{thm:if-q}): a loss that bounds influence on confidently-wrong samples does not over-fit to label noise and instead treats the noise as implicit regularization.

\textbf{(3) The minority-F1 column is contaminated by a metric artifact.} The drop from $0.670$ (FL/GFL clean) to $0.563$ (FL/GFL noisy) appears large but is dominated by a single class. We address this in Section~\ref{sec:exp-f1-tier}.

\subsection{Tier-stratified analysis}
\label{sec:exp-f1-tier}

The minority class set spans three orders of magnitude in training count: from Bot ($1\,363$) and SSH-Patator ($2\,253$) at the upper end to Heartbleed ($8$) and Sql Injection ($15$) at the extreme tail. Reporting unweighted mean F1 across this range conflates two regimes and produces misleading aggregate numbers. We therefore stratify the minority set into two tiers:
\begin{itemize}[leftmargin=1.6em]
\item \textbf{Moderate-rare} (training count $\in [400, 5000]$): Bot, DoS Slowhttptest, DoS slowloris, FTP-Patator, SSH-Patator, Web Attack-Brute Force, Web Attack-XSS.
\item \textbf{Extreme-rare} (training count $< 100$): Heartbleed, Infiltration, Web Attack-SQL Injection.
\end{itemize}
The extreme-rare tier has between $1$ and $7$ test samples per class; per-class F1 is therefore quantized to a small set of values and is statistically unreliable.

\begin{table}[t]
\caption{Mean F1 stratified by class-frequency tier. Moderate-rare: training counts $[400, 5000]$. Extreme-rare: $< 100$ samples.}
\label{tab:f1-tier}
\centering
\small
\setlength{\tabcolsep}{4pt}
\begin{tabular}{l|cc|cc}
\toprule
& \multicolumn{2}{c|}{\textbf{Moderate-rare}} & \multicolumn{2}{c}{\textbf{Extreme-rare}} \\
Loss & $n{=}0.0$ & $n{=}0.20$ & $n{=}0.0$ & $n{=}0.20$ \\
\midrule
CE       & $0.748$ & $0.741$ & $0.139$ & $0.148$ \\
FL       & $0.750$ & $0.740$ & $0.482$ & $0.148$ \\
GCE      & $0.402$ & $0.397$ & $0.000$ & $0.000$ \\
MAE      & $0.241$ & $0.288$ & $0.000$ & $0.000$ \\
GFL      & $0.753$ & $0.740$ & $0.477$ & $0.148$ \\
GFL-$q$  & $0.637$ & $0.743$ & $0.000$ & $0.000$ \\
\bottomrule
\end{tabular}
\end{table}

\begin{figure}[t]
\centering
\includegraphics[width=0.92\linewidth]{figs/fig_f1_tier_curves}
\caption{Moderate-rare tier F1 across noise rates. CE, FL, and GFL track tightly with small degradation. GFL-$q$ shows clear noise-as-regularization, climbing from $0.64$ (clean) to $0.74$ (20\% noise).}
\label{fig:f1-tier}
\end{figure}

Table~\ref{tab:f1-tier} and Figure~\ref{fig:f1-tier} together tell a clean story. On the moderate-rare tier, the degradation slopes ($\Delta = $ noisy $-$ clean) are: CE~$-0.007$, FL~$-0.010$, GFL~$-0.013$, GFL-$q$~$+0.106$. Cross-entropy and focal loss are robust to this noise level on moderate-rare classes; standard GFL is approximately tied with FL. The GFL-$q$ result, however, is a positive +0.106 swing---a direct empirical confirmation of the noise-as-regularization signature predicted by Theorem~\ref{thm:if-q}.

On the extreme-rare tier, FL and GFL show identical $0.482 \to 0.148$ behavior at noise=0, dropping to $0.148$ at noise=0.20. With test-set sample counts of $1$--$7$ per class, every per-class F1 in this tier is quantized to roughly $\{0, 0.5, 1\}$ and a single misclassification swings the tier mean by $\sim 0.10$. We therefore do not draw conclusions from the extreme-rare numbers, and report them only for completeness.

\subsection{Saturation-parameter ablation}
\label{sec:exp-tau}

To test whether the F1-level invariance of GFL to FL persists across saturation strengths, we sweep $\tau \in \{0, 0.3, 0.7, 1.0, 1.5, 2.0\}$ on $\{0, 0.20\}$ noise rates with $3$ seeds.

\begin{table}[t]
\caption{GFL macro-F1 across $\tau$ at $\gamma = 2$. CE and FL baselines listed for reference.}
\label{tab:tau-ablation}
\centering
\small
\begin{tabular}{lcc}
\toprule
Configuration & Clean & $20\%$ noise \\
\midrule
CE              & $0.706 \pm .002$ & $0.704 \pm .001$ \\
FL              & $0.775 \pm .000$ & $0.703 \pm .000$ \\
\midrule
GFL ($\tau{=}0.0$)  & $0.775 \pm .000$ & $0.703 \pm .000$ \\
GFL ($\tau{=}0.3$)  & $0.775 \pm .003$ & $0.704 \pm .000$ \\
GFL ($\tau{=}0.7$)  & $0.774 \pm .000$ & $0.703 \pm .001$ \\
GFL ($\tau{=}1.0$)  & $0.773 \pm .002$ & $0.704 \pm .000$ \\
GFL ($\tau{=}1.5$)  & $0.773 \pm .002$ & $0.703 \pm .001$ \\
GFL ($\tau{=}2.0$)  & $0.774 \pm .002$ & $0.704 \pm .000$ \\
\bottomrule
\end{tabular}
\end{table}

Across the full $\tau$ range $[0, 2]$, GFL macro-F1 varies by less than $0.003$ at either noise level---substantially below the seed-to-seed standard deviation. This is consistent with the empirical observation in Section~\ref{sec:exp-emp}: the test distribution post-training places only $0.4\%$ of samples in the regime where the $\tau$-controlled gradient attenuation has any effect, so the F1 metric is insensitive to $\tau$ on this dataset's natural distribution. The result is a positive finding for the GFL family at the practitioner level: choice of $\tau$ in $[0, 2]$ does not require careful tuning for downstream F1 on CICIDS2017.

% =====================================================================
\section{Discussion}
\label{sec:discussion}

\subsection{Why F1 is dataset-limited on CICIDS2017}
The bounded-influence theorem makes a precise prediction about the per-sample gradient on confidently-wrong samples (Theorem~\ref{thm:grad-comp}c). We confirm that prediction directly to machine precision (Section~\ref{sec:exp-synth}) and to within $10\%$ on real samples (Section~\ref{sec:exp-emp}). Yet the downstream F1 difference between FL and GFL is within seed noise (Tables~\ref{tab:f1-aggregate}, \ref{tab:f1-tier}). The cause is a property of the dataset, not of the loss: on CICIDS2017 with a converged classifier, only $0.4\%$ of samples reach the regime where the bounded-influence advantage applies. This dilutes any per-minibatch effect by a factor of approximately $250$.

This is not a contradiction of the theorem; it is a statement about its applicability domain. Datasets with higher inherent label noise (e.g., naturally corrupted web-scraped image datasets with reported $20$--$40\%$ noise rates), or tasks where confidently-wrong predictions are intrinsic to the difficulty (e.g., long-tail classification on iNaturalist), should produce a larger fraction of samples in the bounded-influence regime and a correspondingly larger F1 effect. Validation in those regimes is left to future work.

\subsection{What did manifest in F1}
The $q$-deformed extension (Theorem~\ref{thm:if-q}) provides uniform B-robustness on the closed interval $[0,1]$, not just at the noise boundary. This is a structurally stronger property than the GFL family achieves and is directly visible in the F1 results: GFL-$q$ exhibits noise-as-regularization on the moderate-rare tier, with mean F1 climbing from $0.637$ (clean) to $0.743$ ($20\%$ noise). This is the cleanest empirical confirmation of the bounded-influence framework in our results.

\subsection{Pilot study reconciled}
The pilot of Section~\ref{sec:pilot} found that focal loss did not improve over cross-entropy when applied to a CNN-LSTM architecture. The current full evaluation finds the same effect at the loss-function level: at default $\tau = 0.3$ on the MLP backbone, FL macro-F1 is $0.775$ versus CE at $0.706$, a $0.069$ gain. The discrepancy between the pilot and the main evaluation is attributable to the architecture: the CNN+LSTM backbone in the pilot was likely already saturating the achievable performance at this dataset scale, leaving no room for the loss-level improvement to manifest. On the simpler MLP backbone, the focal-loss focusing effect on moderate-rare classes does become visible.

% =====================================================================
\section{Limitations and Future Work}
\label{sec:limitations}

We disclose three limitations explicitly. First, the F1-level invariance of FL and GFL on CICIDS2017 (Section~\ref{sec:exp-tau}) is a function of the dataset's $p_t$ distribution and not a refutation of the bounded-influence theorem; cross-dataset validation is required. Second, the extreme-rare tier of CICIDS2017 (Heartbleed, Infiltration, Sql Injection) has too few test samples to support per-class F1 conclusions, and these classes are excluded from our headline numbers. Third, the implementation uses fixed $(\gamma, \tau, q)$ throughout training; an entropy-adaptive schedule for $\tau$ may further improve results and is a natural extension. Future work will validate the framework on CIFAR-10/100 with synthetic label noise (the standard label-noise benchmark, where the bounded-influence regime is more frequently triggered) and on iNaturalist for natural long-tail evaluation.

% =====================================================================
\section{Related Work}
\label{sec:related}

\textbf{Loss-level approaches to imbalance.} Focal loss \cite{lin2017focal} introduced the polynomial modulator that this work generalizes. Class-balanced loss \cite{cui2019class} reweights by effective sample number per class; this is orthogonal and can be applied as the $\alpha_t$ term in our formulation. Asymmetric and varifocal losses for object detection \cite{zhang2021varifocal} modify the modulator for positive vs.\ negative samples but retain the unbounded-gradient property analyzed here.

\textbf{Bounded-influence and noise-robust losses.} Mean Absolute Error (MAE) is uniformly bounded but converges slowly; Generalized Cross-Entropy (GCE) \cite{zhang2018generalized} interpolates between CE and MAE via a power-deformation parameter and is the closest single-parameter analogue to the $q$ component of our $(\gamma, \tau, q)$ family. Symmetric Cross-Entropy and Active-Passive losses extend this line. Our $q$-deformed extension (Definition~\ref{def:gfl-q}) is the first to combine the GCE-style power deformation with a focal modulator and a saturation parameter in a single closed-form construction.

\textbf{Machine learning for network intrusion detection.} The fundamental challenge of applying ML to NIDS---including the gap between benchmark accuracy and operational detection---is analyzed by Sommer and Paxson~\cite{sommer2010outside}, who argue that IDS evaluation must account for the extreme rarity of attack traffic and the cost of false alarms. This motivates tier-stratified evaluation (Section~\ref{sec:exp-f1-tier}) rather than aggregate accuracy.

\textbf{Gradient-norm validation.} Direct measurement of per-sample gradient distributions has been used to characterize loss behavior under noise \cite{ghosh2017robust} and to motivate sample reweighting schemes \cite{ren2018learning}. Our use of bucketed gradient measurements (Section~\ref{sec:exp-emp}) follows this methodology and provides direct evidence for the predicted attenuation factor.

% =====================================================================
\section{Conclusion}

We presented the Generalized Focal Loss, a two-parameter family that fixes the unbounded-influence and focusing--suppression coupling defects of standard focal loss while recovering it as a parametric limit. The bounded-influence theorem is verified directly: to $3 \times 10^{-4}\%$ relative error on synthetic logits, and to within $10\%$ on real CICIDS2017 samples. The $q$-deformed extension achieves uniform B-robustness on the closed interval and exhibits the predicted noise-as-regularization signature on moderate-rare classes ($+0.106$ F1 from clean to $20\%$ label noise). Downstream F1 on CICIDS2017 is dataset-limited by the small fraction of samples reaching the bounded-influence regime ($0.4\%$ post-convergence); we report this honestly and identify cross-dataset validation as the natural next step. The mathematical framework, the closed-form derivatives, the proofs, and the direct gradient-norm validation methodology are the central contribution and are independent of any single dataset's distributional properties.

% =====================================================================
\bibliographystyle{IEEEtran}
\begin{thebibliography}{99}

\bibitem{lin2017focal} T.-Y. Lin, P. Goyal, R. Girshick, K. He, and P. Doll\'ar, ``Focal loss for dense object detection,'' in \emph{ICCV}, 2017, pp.~2980--2988.

\bibitem{chawla2002smote} N.~V. Chawla, K.~W. Bowyer, L.~O. Hall, and W.~P. Kegelmeyer, ``SMOTE: Synthetic minority over-sampling technique,'' \emph{J. Artif. Intell. Res.}, vol.~16, pp.~321--357, 2002.

\bibitem{sharafaldin2018toward} I.~Sharafaldin, A.~H. Lashkari, and A.~A. Ghorbani, ``Toward generating a new intrusion detection dataset and intrusion traffic characterization,'' in \emph{ICISSP}, 2018, pp.~108--116.

\bibitem{huber1981} P.~J. Huber, \emph{Robust Statistics}.\quad Wiley, 1981.

\bibitem{banerjee2005bregman} A.~Banerjee, S.~Merugu, I.~S. Dhillon, and J.~Ghosh, ``Clustering with Bregman divergences,'' \emph{J. Mach. Learn. Res.}, vol.~6, pp.~1705--1749, 2005.

\bibitem{tsallis1988possible} C.~Tsallis, ``Possible generalization of Boltzmann--Gibbs statistics,'' \emph{J. Stat. Phys.}, vol.~52, pp.~479--487, 1988.

\bibitem{ghosh2017robust} A.~Ghosh, H.~Kumar, and P.~S. Sastry, ``Robust loss functions under label noise for deep neural networks,'' in \emph{AAAI}, 2017.

\bibitem{zhang2018generalized} Z.~Zhang and M.~R. Sabuncu, ``Generalized cross entropy loss for training deep neural networks with noisy labels,'' in \emph{NeurIPS}, 2018.

\bibitem{koh2017understanding} P.~W. Koh and P.~Liang, ``Understanding black-box predictions via influence functions,'' in \emph{ICML}, 2017.

\bibitem{cui2019class} Y.~Cui, M.~Jia, T.-Y. Lin, Y.~Song, and S.~Belongie, ``Class-balanced loss based on effective number of samples,'' in \emph{CVPR}, 2019, pp.~9268--9277.

\bibitem{zhang2021varifocal} H.~Zhang, Y.~Wang, F.~Dayoub, and N.~Sunderhauf, ``VarifocalNet: An IoU-aware dense object detector,'' in \emph{CVPR}, 2021.

\bibitem{ren2018learning} M.~Ren, W.~Zeng, B.~Yang, and R.~Urtasun, ``Learning to reweight examples for robust deep learning,'' in \emph{ICML}, 2018.

\bibitem{cherukuri2025deep} A.~K. Cherukuri \emph{et~al.}, ``Intrusion detection system using deep learning for network security,'' \emph{arXiv:2505.05810}, 2025.

\bibitem{zhang2025deep} L.~Zhang \emph{et~al.}, ``A review of deep learning applications in intrusion detection systems,'' \emph{Applied Sciences}, vol.~15, no.~3, p.~1552, 2025.

\bibitem{nesterov2018} Y.~Nesterov, \emph{Lectures on Convex Optimization}, 2nd~ed.\quad Springer, 2018.

\bibitem{kingma2014adam} D.~P. Kingma and J.~Ba, ``Adam: A method for stochastic optimization,'' in \emph{ICLR}, 2015.

\bibitem{lashkari2017cicflowmeter} A.~H. Lashkari, G.~Draper-Gil, M.~S.~I. Mamun, and A.~A. Ghorbani, ``Characterization of encrypted and VPN traffic using time-related features,'' in \emph{ICISSP}, 2017, pp.~407--414.

\bibitem{he2009learning} H.~He and E.~A. Garcia, ``Learning from imbalanced data,'' \emph{IEEE Trans. Knowl. Data Eng.}, vol.~21, no.~9, pp.~1263--1284, 2009.

\bibitem{sommer2010outside} R.~Sommer and V.~Paxson, ``Outside the closed world: On using machine learning for network intrusion detection,'' in \emph{IEEE Symp. Security and Privacy}, 2010, pp.~305--316.

\end{thebibliography}

\end{document}
